% =====================================================================
%  PREAMBULO COMPARTIDO - Documentacion Tecnica de Arquitectura
%  Universidad Nacional de Colombia, Sede Manizales
% =====================================================================

\documentclass[12pt,letterpaper,oneside]{report}

% ------------------------- Codificacion e idioma --------------------
\usepackage[utf8]{inputenc}
\usepackage[T1]{fontenc}
\usepackage[spanish,es-tabla,es-nodecimaldot,es-noshorthands]{babel}
\usepackage{lmodern}
\usepackage{microtype}

% ------------------------- Diseno y geometria --------------------
\usepackage[letterpaper,top=3cm,bottom=3cm,left=3cm,right=2.5cm,headheight=15pt]{geometry}
\usepackage{fancyhdr}
\usepackage{titlesec}
\usepackage{setspace}
\usepackage{parskip}
\onehalfspacing

% ------------------------- Graficos y diagramas --------------------
\usepackage{xcolor}
\usepackage{tikz}
\usetikzlibrary{arrows.meta,positioning,calc,fit,backgrounds,
                decorations.pathreplacing,chains,
                shapes.geometric,shapes.misc,matrix}
\usepackage{float}
\usepackage{graphicx}
\usepackage{adjustbox}

% ------------------------- Tablas --------------------------------
\usepackage{array}
\usepackage{booktabs}
\usepackage{longtable}
\usepackage{tabularx}
\usepackage{multirow}

% ------------------------- Cajas, listados y acronimos --------------
\usepackage{enumitem}
\usepackage{tcolorbox}
\tcbuselibrary{skins,breakable}
\usepackage{listings}
\usepackage[printonlyused]{acronym}
\usepackage{caption}
\usepackage[hidelinks]{hyperref}

% cleveref debe ir despues de hyperref
\usepackage[spanish]{cleveref}

% ------------------------- Colores institucionales ------------------
\definecolor{azulUC}{RGB}{21,101,192}
\definecolor{azulOscuro}{RGB}{13,71,161}
\definecolor{azulClaro}{RGB}{227,242,253}
\definecolor{verdeUC}{RGB}{56,142,60}
\definecolor{verdeClaro}{RGB}{232,245,233}
\definecolor{naranjaUC}{RGB}{239,108,0}
\definecolor{naranjaClaro}{RGB}{255,243,224}
\definecolor{moradoUC}{RGB}{173,20,87}
\definecolor{moradoClaro}{RGB}{243,229,245}
\definecolor{grisUC}{RGB}{96,96,96}
\definecolor{grisClaro}{RGB}{245,245,245}
% FIN PARTE 1

% ------------------------- Hipervinculos ----------------------------
\hypersetup{
    colorlinks=true,
    linkcolor=azulOscuro,
    urlcolor=azulUC,
    citecolor=verdeUC,
    pdftitle={Arquitectura Tecnica del SGDIC - DIC UNC Sede Manizales},
    pdfauthor={Departamento de Informatica y Computacion},
    pdfsubject={Arquitectura de software: frontend, backend Django, base de datos y despliegue},
    bookmarksnumbered=true,
    bookmarksopen=true}

% ------------------------- Encabezados y pies -----------------------
\pagestyle{fancy}
\fancyhf{}
\fancyhead[L]{\small\textcolor{azulOscuro}{\textit{Universidad Nacional de Colombia -- Sede Manizales}}}
\fancyhead[R]{\small\textcolor{azulOscuro}{\thepage}}
\fancyfoot[C]{\small\textcolor{grisUC}{Facultad de Administración -- Departamento de Informática y Computación}}
\renewcommand{\headrulewidth}{0.6pt}
\renewcommand{\footrulewidth}{0.4pt}

% ------------------------- Estilo de capitulos ----------------------
\titleformat{\chapter}[display]
  {\normalfont\huge\bfseries\color{azulOscuro}}
  {\filright\Large\textcolor{azulUC}{\chaptertitlename\ \thechapter}}
  {12pt}
  {\Huge}
  [\vspace{4pt}{\color{azulUC}\titlerule[1.5pt]}]

\titleformat{\section}
  {\normalfont\Large\bfseries\color{azulOscuro}}{\thesection}{1em}{}

\titleformat{\subsection}
  {\normalfont\large\bfseries\color{azulUC}}{\thesubsection}{1em}{}

\titlespacing*{\chapter}{0pt}{-20pt}{30pt}
